\section*{Abstract}
%A long-standing grand challenge in computer science, accelerating and automating software engineering, has entered a new era with the advent of Large Language Models (LLMs) agents. Coding agents have quickly become a mainstay in professional software development, and therefore require evaluations that represents the complexity of enterprise software engineering. 
We introduce \benchmarkname, a substantially more challenging benchmark that builds upon the best practices of SWE-Bench~\cite{zhang2025swebench}, but is explicitly designed to capture realistic, complex, enterprise-level problems beyond the scope of SWE-Bench.
\benchmarkname~contains 1,865 problems sourced from a diverse set of 41 actively maintained repositories spanning business applications, B2B services, and developer tools. The benchmark is partitioned into a \emph{public} set with open access to problems sourced from 11 repositories, a \emph{held-out} set of 12 repositories and a \emph{commercial} set of 18 proprietary repositories where we have formal partnership agreements with early-stage startups. Problems in the held-out and the commercial set are not publicly accessible, but we release results on the commercial set.
% 18 of these repositories are proprietary codebases sourced through formal partnership agreements with early-stage startups, forming an additional set of difficult, commercial-grade environments. Our efforts result in a large-scale agent benchmark with 1865 instances (731 public instances) across 41 repositories (11 in the public set, 12 in the held-out set, and 18 in the co). 
%We exclusively use repositories with strong copy-left licenses (e.g., GPL), along with commercial, proprietary codebases sourced from startups through partnership agreements to ensure non-contamination by design. 
% To further mitigate data leakage and enforce stricter generalization, we additionally secured 18 proprietary commercial codebases from early-stage startups through formal partnership agreements.
Our benchmark features long-horizon tasks that may require hours to days for a professional software engineer to complete, often involving patches across multiple files and substantial code modifications. All tasks are human-verified and augmented with sufficient context to ensure resolvability. In our evaluation of widely used coding models, under a unified scaffold, we observe that their performance on \benchmarkname~remains below 45\% (Pass@1). To better understand these limitations, we cluster the failure modes observed in the collected agent trajectories for a clearer characterization of the error patterns exhibited by current models. Overall, \benchmarkname~provides a contamination-resistant testbed that more faithfully captures the complexity and diversity of real-world software development, advancing the pursuit of truly autonomous software engineering agents at a professional level. 


% A long-standing grand challenge in computer science, accelerating and automating software engineering, has entered a new era with the advent of Large Language Models (LLMs). Coding agents have quickly become a mainstay in professional software development. Despite surpassing top human performers in competitive programming and achieving over 70\% accuracy on benchmarks such as SWE-Bench, many practitioners report conflicted feedback on their real-world effectiveness. Does this really mean that today's coding agents are already capable of solving 70\% of software engineering tasks? We argue that existing benchmarks fall short of capturing the full complexity of industrial challenges. \sail{This sentence "We argue ... industrial challenges" is a bit annoying. It gave an argument but did not mention why. I suggest to add another sentence describing the main claims "why they fall short". } We introduce \benchmarkname, a novel benchmark designed from the ground up to provide a more rigorous and realistic evaluation of AI agents for software engineering. \benchmarkname\ addresses key shortcomings in previous benchmarks through three foundational principles: (1) diverse task selection from complex, actively maintained repositories spanning consumer applications, B2B services, and developer tools, better representing industrial software engineering challenges; (2) difficult, human-verified tasks selected to be representative of professional software engineering tasks, with patches spanning across multiple files and requiring a significant code change, and (3) exclusive use of repositories with strong copy-left licenses (e.g., GPL) and commercial, proprietary codebases sourced from startups to ensure non-contamination by design.  Through a comprehensive evaluation of state-of-the-art models on \benchmarkname, we establish new performance baselines and provide critical insights into current model capabilities and limitations. Top models, such as Opus 4.1 and GPT-5, score 23\% (Pass@1) on \benchmarkname, compared to 70\%+ on SWE-Bench Verified. Our benchmark offers a  contamination-resistant testbed that more accurately reflects the complexity and diversity of real-world software development, advancing the pursuit of truly autonomous professional-level software engineering agents.

% \sail{Overall I think the abstract is clear and straightforward. My general comment is to replace hand-wavy terms with accurate terms. The three contributions of our benchmark needs some sorting -- (3) sounds more important and the rest because that is the most important thing we did here. Can we talk a bit more in the abstract how we acquire these start-up code base. This point is actually important but is not well-discussed in the intro. Please talk about the number of tasks and the diversity distribution in (1). We at least need to tell people how many tasks we have. Another thing is currently missing from the abstract is that: what eval process we follow. Is that same as SWE-Bench?  }